% !TEX TS-program = pdflatex
% !TeX program = pdflatex
% !TEX encoding = UTF-8
% !TEX spellcheck = en

\documentclass[xcolor=table]{beamer}

\usepackage{../extra/beamer/karimml}

\input{options2}

\subtitle[CNNs]{Convolutional Neural Networks}  

\changegraphpath{../img/CNN/}

\begin{document}

\begin{frame}
	\frametitle{\inserttitle}
	\framesubtitle{\insertshortsubtitle: Motivation}


\end{frame}

\begin{frame}
	\frametitle{\inserttitle}
	\framesubtitle{\insertshortsubtitle: Plan}
	
	\begin{multicols}{2}
		%	\small
		\tableofcontents
	\end{multicols}
\end{frame}

\begin{frame}
	\frametitle{\inserttitle}
	\framesubtitle{\insertshortsubtitle: A little fun}
	
	\hgraphpage[\textwidth]{humor/humor-convnet.jpg}
	
\end{frame}

\section{Traditional image processing}

\begin{frame}
	
	\frametitle{\insertshortsubtitle}
	\framesubtitle{\insertsection: Architecture}
	
	\begin{center}
		\hgraphpage[\textwidth]{img-learn.pdf}
	\end{center}

\end{frame}

\begin{frame}
	\frametitle{\insertshortsubtitle: \insertsection}
	\framesubtitle{\insertsubsection: Convolution}
	
	\begin{itemize}
		\item Many preprocessing techniques exist (see \cite{2010-alginahi})
		\item Here, we are interested in convolution (modify the value of a pixel based on its neighbors)
		\item Two parameters: \optword{padding} (Surround the image with 0s in order to preserve the original size), \optword{stride} (the kernel/mask's scrolling step).
	\end{itemize}
	
	\begin{center}
		\hgraphpage[.8\textwidth]{conv.pdf}
	\end{center}

\end{frame}

\begin{frame}
	\frametitle{\insertshortsubtitle: \insertsection}
	\framesubtitle{\insertsubsection: Example of kernels}
	
	Ua example from \href{https://fr.wikipedia.org/wiki/Noyau\_(traitement\_d\%27image)}{Wikipedia}
	
	\begin{tabular}{p{.3\textwidth}p{.1\textwidth}p{.2\textwidth}p{.2\textwidth}}
		\hline\hline
		Operation & original & Kernel & Transformed \\
		\hline
		Contour detection & 
		\graphpage[valign=c]{Vd-Orig.png} & 
		$\begin{bmatrix}
		-1 & -1 & -1\\ 
		-1 & 8 & -1\\ 
		-1 & -1 & -1
		\end{bmatrix}$ & 
		\graphpage[valign=c]{Vd-Edge3.png} \\
		
		\hline
		Sharpness improvement & 
		\graphpage[valign=c]{Vd-Orig.png} & 
		$\begin{bmatrix}
		0 & -1 & 0\\ 
		-1 & 5 & -1\\ 
		0 & -1 & 0
		\end{bmatrix}$ & 
		\graphpage[valign=c]{Vd-Sharp.png} \\
		
		\hline
		Box blur & 
		\graphpage[valign=c]{Vd-Orig.png} & 
		$\frac{1}{9}\begin{bmatrix}
		1 & 1 & 1\\ 
		1 & 1 & 1\\ 
		1 & 1 & 1
		\end{bmatrix}$ & 
		\graphpage[valign=c]{Vd-Blur2.png} \\
		\hline\hline
		
	\end{tabular}

\end{frame}

\begin{frame}
	\frametitle{\insertshortsubtitle: \insertsection}
	\framesubtitle{\insertsubsection: Limitations}
	
	According to LeCun and his colleagues \cite{1998-lecun}:
	\begin{itemize}
		\item When the size of the images is large, there will be a large number of parameters to train.
		\item To achieve this, a large dataset must be provided.
		\item The memory required to store the parameters will be very large.
		\item To train the model on images, preprocessing such as translations and distortions must be applied.
		\item Variations in images (such as the position of an object) can only be captured when multiple layers are used.
	\end{itemize}

\end{frame}

\section{Convolutional layers}

\begin{frame}
	\frametitle{\insertshortsubtitle}
	\framesubtitle{\insertsection}
	
\end{frame}

\begin{frame}
	\frametitle{\insertshortsubtitle: \insertsection}
	\framesubtitle{\insertsubsection: Conv2D layer}

	\begin{minipage}{0.60\textwidth} 
		\begin{itemize}
			\item Preserve images' spatial structure.
			\item $ w' = \frac{w - w_f + 2P}{S} + 1$, $ h' = \frac{h - h_f + 2P}{S} + 1$
			\item A layer can have $k$ filters/kernels.
			\item Number of parameters: $w_f * h_f * c * k$ plus $k$ biases.
			\item Example, \expword{image: 32x32x3;} \expword{kernel: 5x5; s: 1; p: 0; k: 6.} \\ Number of trainable parameters: \expword{456}
		\end{itemize}
	\end{minipage}
	%
	\begin{minipage}{0.39\textwidth}
		\hgraphpage[\textwidth]{conv2d.pdf}
		
		\hgraphpage[.7\textwidth]{conv2d-exp_.pdf}
	\end{minipage}

\end{frame}

\begin{frame}
	\frametitle{\insertshortsubtitle: \insertsection}
	\framesubtitle{\insertsubsection: Conv2D layer (Example)}
	
	
	\begin{center}
		\vskip-6pt\hgraphpage[0.8\textwidth]{conv2D_imlp.pdf}
	\end{center}\vskip-16pt
	
	{\scriptsize 
		\[O_{11} = ReLU(\Sigma = b + F_{11} I_{11} + \textcolor{blue}{F_{12} I_{12}} + F_{21} I_{21} + \textcolor{red}{F_{22} I_{22}})  \]
		\[O_{12} = ReLU(\Sigma = b + \textcolor{blue}{F_{11} I_{12}} + F_{12} I_{13} + \textcolor{red}{F_{21} I_{22}} + F_{22} I_{23})  \]
		\[O_{21} = ReLU(\Sigma = b + F_{11} I_{21} + \textcolor{red}{F_{12} I_{22}} + F_{21} I_{31} + F_{22} I_{32})  \]
		\[O_{22} = ReLU(\Sigma = b + \textcolor{red}{F_{11} I_{22}} + F_{12} I_{23} + F_{21} I_{32} + F_{22} I_{33})  \]
		
		\[\textcolor{red}{\frac{\partial J}{\partial I_{22}}} 
		= \underbrace{\frac{\partial J}{\partial O_{11}}}_{G_{11} = 1} 
		\underbrace{\frac{\partial ReLU}{\partial \sum}}_{1} 
		\underbrace{\frac{\partial \sum}{\partial I_{22}}}_{F_{22} = 3}
		+ \underbrace{\frac{\partial J}{\partial O_{12}}}_{G_{12} = 2} 
		\underbrace{\frac{\partial ReLU}{\partial \sum}}_{1} 
		\underbrace{\frac{\partial \sum}{\partial I_{22}}}_{F_{21} = -2}
		+ \underbrace{\frac{\partial J}{\partial O_{21}}}_{G_{21} = 3} 
		\underbrace{\frac{\partial ReLU}{\partial \sum}}_{1} 
		\underbrace{\frac{\partial \sum}{\partial I_{22}}}_{F_{12} = -1}
		+ \underbrace{\frac{\partial J}{\partial O_{22}}}_{G_{22} = 4} 
		\underbrace{\frac{\partial ReLU}{\partial \sum}}_{0} 
		\underbrace{\frac{\partial \sum}{\partial I_{22}}}_{F_{11} = 1}\]
		
		\[\textcolor{blue}{\frac{\partial J}{\partial I_{12}}} 
		= \underbrace{\frac{\partial J}{\partial O_{11}}}_{G_{11} = 1} 
		\underbrace{\frac{\partial ReLU}{\partial \sum}}_{1} 
		\underbrace{\frac{\partial \sum}{\partial I_{12}}}_{F_{12} = -1}
		+ \underbrace{\frac{\partial J}{\partial O_{12}}}_{G_{12} = 2} 
		\underbrace{\frac{\partial ReLU}{\partial \sum}}_{1} 
		\underbrace{\frac{\partial \sum}{\partial I_{12}}}_{F_{11} = 1}\]
	}
	
\end{frame}

\section{Pooling layers}

\begin{frame}
	\frametitle{\insertshortsubtitle}
	\framesubtitle{\insertsection}
	
\end{frame}

\begin{frame}
\frametitle{Convolutional neural networks}
\framesubtitle{Pooling layers}

\begin{minipage}{0.60\textwidth} 
	\begin{itemize}
		\item Makes the representation smaller and more manageable.
		\item $ w' = \frac{w - w_f + 2P}{S} + 1$,  $ h' = \frac{h - h_f + 2P}{S} + 1$
		\item No parameters to train.
		\item \optword{Max Pool}: the gradient go back to the wining cell (the max).
		\item \optword{Average Pool}: the gradient passed to the participated cells is $\frac{\text{gradient}}{w_f * h_f}$. 
	\end{itemize}
\end{minipage}
%
\begin{minipage}{0.39\textwidth}
	\hgraphpage[\textwidth]{maxpool.pdf}
	
	\hgraphpage[.8\textwidth]{pool-exp_.pdf}
\end{minipage}

\end{frame}


\end{document}